%%%%%%%%%%%%%%%%%%%%%%%%%%%%%%%%%%%%%%%%%%%%%%%%%%%%%%%%%%%%%%%%%%%%%%%%%%%%%%%
% TURBOQUANT and online vector quantization
%%%%%%%%%%%%%%%%%%%%%%%%%%%%%%%%%%%%%%%%%%%%%%%%%%%%%%%%%%%%%%%%%%%%%%%%%%%%%%%

% Fallback for standalone compilation if mathtools is not loaded.
\providecommand{\coloneqq}{\mathrel{\mathop:}=}

\subsection{TURBOQUANT: Online Vector Quantization via Random Rotations}
\label{subsec:turboquant}

In the previous parts of the course, we studied the operational and information-theoretic rate--distortion functions, the Gaussian rate--distortion function, and the idea that a random rotation can make an arbitrary high-dimensional vector look ``Gaussian'' in its coordinates.
TURBOQUANT is a useful modern example where these ideas directly lead to a practical quantization algorithm.
The main message is:
\begin{center}
\emph{Randomly rotate a vector, quantize each coordinate with a scalar quantizer matched to the rotated-coordinate distribution, and rotate back.}
\end{center}
For mean-squared error, this already gives a nearly rate--distortion optimal vector quantizer.
For inner products, one more ingredient is needed: an unbiased one-bit quantizer for the residual.

\paragraph{Why this topic is different from ordinary scalar quantization.}
Scalar quantization takes a coordinate $x_j$ and rounds it to one of a small number of levels.
This is simple, but it is usually not optimal for high-dimensional vectors because coordinates can be highly nonuniform: a few coordinates can be very large, while most are small.
For example, in LLM KV caches or embedding vectors, coordinate outliers can dominate low-bit scalar quantization.
A random rotation spreads the energy across coordinates, so that after rotation every coordinate has a predictable distribution and typical magnitude about $1/\sqrt d$.
This turns a worst-case vector quantization problem into a nearly scalar problem with a known source law.

\paragraph{Relation to Gaussian rate--distortion.}
Recall that if $X\sim \mathcal N(0,\sigma^2)$ and squared-error distortion is used, then the Gaussian distortion--rate function in bits is
\begin{equation}
    D(R)=\sigma^2 2^{-2R}.
\end{equation}
If a $d$-dimensional vector has unit energy, each rotated coordinate has variance roughly $1/d$.
Quantizing each coordinate at $b$ bits suggests a total MSE scaling of
\begin{equation}
    d\cdot \frac1d 2^{-2b}=4^{-b}.
\end{equation}
TURBOQUANT achieves this same exponential dependence on bit-width $b$, up to a small constant factor, without assuming that the input vectors are drawn from a Gaussian distribution.

\subsubsection*{Problem formulation}

Let $x\in \mathbb R^d$ be a vector to be compressed.
A vector quantizer consists of two maps
\begin{equation}
    Q: \mathbb R^d\to \{0,1\}^{B},
    \qquad
    Q^{-1}:\{0,1\}^{B}\to \mathbb R^d,
\end{equation}
where $B$ is the total number of bits.
We write $B=bd$, where $b$ is the \emph{bit-width}, i.e., bits per coordinate.
The reconstructed vector is
\begin{equation}
    \widetilde x = Q^{-1}(Q(x)).
\end{equation}

Most of the theory is stated for unit vectors $x\in S^{d-1}$, where
\begin{equation}
    S^{d-1}=\{x\in \mathbb R^d:\|x\|_2=1\}.
\end{equation}
This is not a serious restriction.
For a general nonzero vector $x$, write
\begin{equation}
    x=\|x\|_2 u,
    \qquad u\in S^{d-1}.
\end{equation}
We quantize $u$ and store the scalar norm $\|x\|_2$ as side information.
For large $d$, the cost of one scalar norm per vector is negligible on a per-coordinate basis.
The same remark applies later to the residual norm stored by the inner-product version.

\paragraph{MSE distortion.}
For reconstruction error, define
\begin{equation}
    D_{\mathrm{mse}}(Q;x)
    \coloneqq
    \mathbb E_Q\left[\|x-Q^{-1}(Q(x))\|_2^2\right],
\end{equation}
where the expectation is over the randomness of the quantizer, such as the random rotation.
For unit vectors, the natural target is a distortion on the order of $4^{-b}$.

\paragraph{Inner-product distortion.}
In many applications, the vector is not used by itself.
It is used through inner products.
For example, attention uses inner products between queries and keys, and nearest-neighbor search ranks database vectors by inner product or cosine similarity with a query.
For a query vector $y\in\mathbb R^d$, define
\begin{equation}
    D_{\mathrm{prod}}(Q;x,y)
    \coloneqq
    \mathbb E_Q\left[
        \left|\langle y,x\rangle-
        \langle y,Q^{-1}(Q(x))\rangle\right|^2
    \right].
\end{equation}
For inner products, low MSE is not the only desirable property.
We often also want the estimator to be unbiased:
\begin{equation}
    \mathbb E_Q\left[\langle y,Q^{-1}(Q(x))\rangle\right]
    =\langle y,x\rangle.
\end{equation}
The MSE-optimized version of TURBOQUANT is generally biased for inner products.
The inner-product version fixes this by quantizing the residual with a one-bit unbiased estimator.

\subsubsection*{The random rotation idea}

Let $\Pi\in\mathbb R^{d\times d}$ be a uniformly random orthogonal matrix, also called a Haar random rotation.
For a fixed unit vector $x\in S^{d-1}$, define
\begin{equation}
    Y=\Pi x.
\end{equation}
Because $\Pi$ is uniformly random, $Y$ is uniformly distributed on the unit sphere $S^{d-1}$, no matter what $x$ was.
Thus, the algorithm does not need to know the distribution of the input data.
The random rotation itself creates a known coordinate distribution.

\begin{lemma}[One coordinate of a random point on the sphere]
\label{lem:tq-sphere-coordinate}
Let $Y\sim\mathrm{Unif}(S^{d-1})$.
Then each coordinate $Y_j$ has density
\begin{equation}
    f_d(t)
    =
    \frac{\Gamma(d/2)}{\sqrt\pi\,\Gamma((d-1)/2)}
    (1-t^2)^{(d-3)/2},
    \qquad -1\le t\le 1.
\end{equation}
Equivalently, $Y_j^2$ follows a Beta distribution with parameters $1/2$ and $(d-1)/2$.
For large $d$,
\begin{equation}
    \sqrt d\,Y_j \Rightarrow \mathcal N(0,1).
\end{equation}
\end{lemma}

\paragraph{Proof idea.}
Fix the first coordinate $Y_1=t$.
The remaining coordinates lie on a $(d-2)$-dimensional sphere of radius $\sqrt{1-t^2}$.
The surface area of that slice is proportional to $(1-t^2)^{(d-3)/2}$.
Normalizing over $t\in[-1,1]$ gives the density above.
The Gaussian approximation follows by writing
\begin{equation}
    (1-t^2)^{(d-3)/2}
    \approx \exp\left(-\frac{d t^2}{2}\right)
\end{equation}
when $t$ is of order $1/\sqrt d$.

\paragraph{Near independence.}
The coordinates of $Y$ are not exactly independent because
\begin{equation}
    \sum_{j=1}^{d}Y_j^2=1.
\end{equation}
However, any fixed number of coordinates becomes asymptotically independent after scaling by $\sqrt d$.
More precisely, for fixed $k\ll d$,
\begin{equation}
    \sqrt d\,(Y_1,\dots,Y_k)
    \Rightarrow
    \mathcal N(0,I_k).
\end{equation}
This is the key heuristic that justifies scalar quantization after rotation: each coordinate behaves almost like an independent sample from $\mathcal N(0,1/d)$.

\paragraph{A small numerical picture.}
For a unit vector in dimension $d=1000$, a typical rotated coordinate has standard deviation $1/\sqrt{1000}\approx 0.0316$.
A one-bit Gaussian Lloyd--Max quantizer uses two levels
\begin{equation}
    \pm \sqrt{\frac{2}{\pi d}}
    \approx \pm 0.0252.
\end{equation}
Thus the quantizer is not trying to represent arbitrary real numbers.
After rotation, it only needs to represent many small coordinates whose distribution is tightly concentrated near zero.

\begin{center}
\begin{tabular}{c|c|c|c}
Original vector & Random rotation & Scalar quantize & Rotate back \\
\hline
$x$ & $y=\Pi x$ & $\widetilde y_j\in\{c_1,\dots,c_{2^b}\}$ & $\widetilde x=\Pi^\top \widetilde y$
\end{tabular}
\end{center}

\subsubsection*{MSE-optimized TURBOQUANT}

The MSE version of TURBOQUANT uses a scalar quantizer matched to the density $f_d$ in Lemma~\ref{lem:tq-sphere-coordinate}.
Let $K=2^b$ be the number of scalar reconstruction levels.
Choose ordered centroids
\begin{equation}
    -1\le c_1\le c_2\le \cdots\le c_K\le 1.
\end{equation}
With squared-error distortion, the decision boundary between $c_i$ and $c_{i+1}$ is the midpoint
\begin{equation}
    \tau_i=\frac{c_i+c_{i+1}}2.
\end{equation}
Set $\tau_0=-1$ and $\tau_K=1$.
The optimal scalar cost is the continuous one-dimensional $K$-means cost
\begin{equation}
\label{eq:tq-scalar-cost}
    C(f_d,b)
    \coloneqq
    \min_{c_1\le\cdots\le c_K}
    \sum_{i=1}^{K}
    \int_{\tau_{i-1}}^{\tau_i} (t-c_i)^2 f_d(t)\,dt.
\end{equation}
This is a Lloyd--Max problem.
The codebook can be computed once for each $d$ and $b$, stored, and reused for all vectors.
No dataset-specific codebook training is required.

\paragraph{Algorithm: \textsc{TURBOQUANT}$_{\mathrm{mse}}$.}
For a unit vector $x\in S^{d-1}$:
\begin{enumerate}[leftmargin=2em]
    \item Draw or fix a public random orthogonal matrix $\Pi$.
    \item Rotate the vector:
    \[
        y=\Pi x.
    \]
    \item For each coordinate $j$, find the closest centroid:
    \[
        \mathrm{idx}_j=\arg\min_{i\in\{1,\dots,2^b\}} |y_j-c_i|.
    \]
    Store the $d$ indices, each using $b$ bits.
    \item To dequantize, set
    \[
        \widetilde y_j=c_{\mathrm{idx}_j}
    \]
    and rotate back:
    \[
        \widetilde x=\Pi^\top \widetilde y.
    \]
\end{enumerate}

\paragraph{Low-bit examples.}
In moderate or high dimension, $f_d$ is very close to $\mathcal N(0,1/d)$.
For this Gaussian approximation:
\begin{center}
\begin{tabular}{c|c|c}
Bit-width $b$ & Approximate scalar centroids & Total MSE \\
\hline
$1$ & $\left\{-\sqrt{2/(\pi d)},\sqrt{2/(\pi d)}\right\}$ & $1-2/\pi\approx 0.363$ \\
$2$ & $\{\pm 0.453/\sqrt d,\ \pm 1.51/\sqrt d\}$ & $\approx 0.117$ \\
$3$ & Lloyd--Max levels for $\mathcal N(0,1/d)$ & $\approx 0.03$ \\
$4$ & Lloyd--Max levels for $\mathcal N(0,1/d)$ & $\approx 0.009$
\end{tabular}
\end{center}
The one-bit value $1-2/\pi$ is worth remembering.
It is exactly the MSE of the two-level quantizer $\widehat Y=\sqrt{2/(\pi d)}\operatorname{sign}(Y)$ when $Y\sim\mathcal N(0,1/d)$.
Indeed,
\begin{align}
    \mathbb E[(Y-a\operatorname{sign}Y)^2]
    &=\mathbb E[Y^2]-2a\mathbb E|Y|+a^2,\\
    a^\star&=\mathbb E|Y|=\sqrt{\frac{2}{\pi d}},\\
    d\cdot \mathbb E[(Y-a^\star\operatorname{sign}Y)^2]
    &=1-\frac{2}{\pi}.
\end{align}

\begin{theorem}[MSE guarantee for \textsc{TURBOQUANT}$_{\mathrm{mse}}$]
\label{thm:tq-mse}
For every unit vector $x\in S^{d-1}$, the MSE-optimized TURBOQUANT reconstruction satisfies
\begin{equation}
    D_{\mathrm{mse}}(Q_{\mathrm{mse}};x)
    =\mathbb E\|x-\widetilde x\|_2^2
    \le
    \frac{\sqrt 3\,\pi}{2}\,4^{-b}.
\end{equation}
For small bit-widths, a sharper numerical evaluation gives approximately
\begin{equation}
    D_{\mathrm{mse}}\approx
    0.36,
    \quad 0.117,
    \quad 0.03,
    \quad 0.009
\end{equation}
for $b=1,2,3,4$, respectively.
\end{theorem}

\paragraph{Proof.}
Let $Y=\Pi x$ and let $\widetilde Y$ be the coordinatewise scalar reconstruction.
Since $\Pi$ is orthogonal,
\begin{equation}
    \|x-\widetilde x\|_2^2
    =\|\Pi x-\widetilde Y\|_2^2
    =\|Y-\widetilde Y\|_2^2.
\end{equation}
Therefore,
\begin{align}
    D_{\mathrm{mse}}
    &=\mathbb E\sum_{j=1}^{d}(Y_j-\widetilde Y_j)^2\\
    &=\sum_{j=1}^{d}\mathbb E(Y_j-c_{\mathrm{idx}_j})^2\\
    &=d\,C(f_d,b),
\end{align}
where the last equality uses the fact that all coordinates have the same marginal density $f_d$ and the same scalar Lloyd--Max rule is applied to each coordinate.
For larger $b$, the Panter--Dite high-resolution scalar quantization formula gives
\begin{equation}
    C(f_d,b)
    \le
    \frac1{12}
    \left(\int f_d(t)^{1/3}\,dt\right)^3 4^{-b}
    =
    \frac{\sqrt 3\,\pi}{2d}\,4^{-b}.
\end{equation}
Multiplying by $d$ gives the theorem.
\hfill$\square$

\paragraph{Interpretation.}
The proof is short because the random rotation did the hard work.
It changed an arbitrary unit vector into a random point on the sphere.
The scalar quantizer sees a known coordinate distribution, and the total distortion is just $d$ times the scalar distortion.
The result is a vector quantizer whose distortion scales like $4^{-b}$, matching the Gaussian rate--distortion exponent.

\subsubsection*{Information-theoretic lower bound}

We now explain why the $4^{-b}$ scaling cannot be improved in general.
This is the rate--distortion part of the story.

\begin{lemma}[Shannon lower bound for MSE]
\label{lem:tq-slb}
Let $X\in\mathbb R^d$ have differential entropy $h(X)$ measured in bits.
For any reconstruction $\widehat X$ satisfying
\begin{equation}
    I(X;\widehat X)\le B
\end{equation}
and
\begin{equation}
    \mathbb E\|X-\widehat X\|_2^2\le D,
\end{equation}
we have
\begin{equation}
    D\ge
    \frac{d}{2\pi e}
    2^{\frac{2}{d}(h(X)-B)}.
\end{equation}
\end{lemma}

\paragraph{Proof.}
Let $Z=X-\widehat X$.
Then
\begin{align}
    B
    &\ge I(X;\widehat X)\\
    &=h(X)-h(X\mid \widehat X)\\
    &=h(X)-h(Z\mid \widehat X)\\
    &\ge h(X)-h(Z).
\end{align}
Among all $d$-dimensional random vectors with second moment $\mathbb E\|Z\|_2^2\le D$, the Gaussian with covariance $(D/d)I_d$ maximizes entropy.
Thus
\begin{equation}
    h(Z)
    \le \frac d2\log_2\left(2\pi e\frac{D}{d}\right).
\end{equation}
Rearranging gives the stated lower bound.
\hfill$\square$

\paragraph{Lower bound on the sphere.}
A version of the Shannon lower bound for the uniform distribution on $S^{d-1}$ gives
\begin{equation}
\label{eq:tq-sphere-lower}
    D(B)\ge 2^{-2B/d}.
\end{equation}
The quick way to remember this is to use the surface area
\begin{equation}
    A_d=\frac{2\pi^{d/2}}{\Gamma(d/2)}
\end{equation}
and the surface entropy $h(X)=\log_2 A_d$ for $X\sim\mathrm{Unif}(S^{d-1})$.
Stirling's approximation gives
\begin{equation}
    A_d^{2/d}\approx \frac{2\pi e}{d},
\end{equation}
so Lemma~\ref{lem:tq-slb} reduces to \eqref{eq:tq-sphere-lower}.
If $B=bd$, then
\begin{equation}
    D(B)\ge 2^{-2b}=4^{-b}.
\end{equation}

\paragraph{From average-case to worst-case.}
The lower bound above is for a random vector uniformly distributed on the sphere.
To turn it into a worst-case lower bound for randomized algorithms, one uses Yao's minimax principle:
\begin{center}
\emph{If every deterministic quantizer has large average distortion for a certain input distribution, then every randomized quantizer has large distortion for at least one input.}
\end{center}
Therefore, for every randomized $b$-bit-per-coordinate quantizer, there exists a hard unit vector $x$ such that
\begin{equation}
    \mathbb E\|x-Q^{-1}(Q(x))\|_2^2\ge 4^{-b}.
\end{equation}
Combining this lower bound with Theorem~\ref{thm:tq-mse}, TURBOQUANT is within a constant factor of the best possible worst-case distortion:
\begin{equation}
    4^{-b}
    \le
    D_{\mathrm{mse}}^{\star}(b)
    \le
    D_{\mathrm{mse}}(Q_{\mathrm{mse}})
    \le
    \frac{\sqrt 3\,\pi}{2}4^{-b}.
\end{equation}
The constant $\frac{\sqrt 3\pi}{2}\approx 2.72$ is not the main point.
The important point is that the exponent $4^{-b}$ is optimal.

\subsubsection*{Why MSE quantization is biased for inner products}

MSE measures Euclidean reconstruction error.
It does not automatically imply that inner products are estimated without bias.
The issue is easiest to see at one bit.
Under the Gaussian approximation, the one-bit MSE quantizer in the rotated domain is
\begin{equation}
    \widetilde Y_j=a\operatorname{sign}(Y_j),
    \qquad
    a=\sqrt{\frac{2}{\pi d}}.
\end{equation}
The reconstructed vector is
\begin{equation}
    \widetilde x=\Pi^\top \widetilde Y.
\end{equation}
By rotational symmetry,
\begin{equation}
    \mathbb E[\widetilde x\mid x]=\alpha x
\end{equation}
for some scalar $\alpha$.
To find $\alpha$, take the inner product with $x$:
\begin{align}
    \alpha
    &=\mathbb E\langle x,\widetilde x\rangle\\
    &=\mathbb E\langle \Pi x,\widetilde Y\rangle\\
    &=\mathbb E\sum_{j=1}^{d} a|Y_j|\\
    &=d\,a\,\mathbb E|Y_1|.
\end{align}
Using $Y_1\approx\mathcal N(0,1/d)$,
\begin{equation}
    a=\mathbb E|Y_1|=\sqrt{\frac{2}{\pi d}},
\end{equation}
so
\begin{equation}
    \alpha\approx d\cdot \frac{2}{\pi d}=\frac{2}{\pi}.
\end{equation}
Hence
\begin{equation}
    \mathbb E\langle y,\widetilde x\rangle
    \approx
    \frac{2}{\pi}\langle y,x\rangle.
\end{equation}
This is a multiplicative shrinkage bias.
The MSE reconstruction is close to $x$, but its average projection in the direction of $x$ is too small.
For ranking, attention, and nearest-neighbor search, this bias can be problematic.

\paragraph{Geometric reason for the bias.}
The one-bit MSE reconstruction has approximately constant magnitude in every rotated coordinate.
Its squared norm is about
\begin{equation}
    \|\widetilde Y\|_2^2=d\cdot \frac{2}{\pi d}=\frac{2}{\pi}.
\end{equation}
Thus the reconstruction has norm about $\sqrt{2/\pi}<1$.
It is a good MSE representative of a sign cell, but it is not an unbiased representation of the original vector for linear measurements.

\subsubsection*{QJL: a one-bit unbiased inner-product quantizer}

The inner-product version of TURBOQUANT uses a one-bit Quantized Johnson--Lindenstrauss transform, abbreviated QJL.
The role of QJL is to represent a residual vector $r$ with one bit per coordinate while preserving inner products in expectation.

Let $S\in\mathbb R^{d\times d}$ have i.i.d. entries $S_{ij}\sim\mathcal N(0,1)$.
For a unit vector $u\in S^{d-1}$, define
\begin{equation}
    Q_{\mathrm{qjl}}(u)=\operatorname{sign}(Su)\in\{-1,+1\}^{d}.
\end{equation}
For a sign vector $z\in\{-1,+1\}^d$, define
\begin{equation}
    Q_{\mathrm{qjl}}^{-1}(z)
    =\frac{\sqrt{\pi/2}}{d}S^\top z.
\end{equation}

\begin{lemma}[QJL unbiasedness and variance]
\label{lem:tq-qjl}
For $u\in S^{d-1}$ and $y\in\mathbb R^d$,
\begin{equation}
    \mathbb E_S\left[\left\langle y,Q_{\mathrm{qjl}}^{-1}(Q_{\mathrm{qjl}}(u))\right\rangle\right]
    =\langle y,u\rangle.
\end{equation}
Moreover,
\begin{equation}
    \operatorname{Var}_S\left(
    \left\langle y,Q_{\mathrm{qjl}}^{-1}(Q_{\mathrm{qjl}}(u))\right\rangle
    \right)
    \le \frac{\pi}{2d}\|y\|_2^2.
\end{equation}
\end{lemma}

\paragraph{Proof.}
Let $s_i^\top$ be the $i$-th row of $S$.
Then
\begin{equation}
    \left\langle y,Q_{\mathrm{qjl}}^{-1}(Q_{\mathrm{qjl}}(u))\right\rangle
    =\frac1d\sum_{i=1}^{d}\sqrt{\frac\pi2}
    (s_i^\top y)\operatorname{sign}(s_i^\top u).
\end{equation}
For a standard Gaussian vector $s$ and a unit vector $u$,
\begin{equation}
    \mathbb E[s\operatorname{sign}(s^\top u)]
    =\sqrt{\frac2\pi}u.
\end{equation}
Therefore each summand has expectation
\begin{equation}
    \sqrt{\frac\pi2}\left\langle y,\sqrt{\frac2\pi}u\right\rangle
    =\langle y,u\rangle,
\end{equation}
and averaging $d$ independent copies keeps the same expectation.
For the variance, define
\begin{equation}
    Z_i=\sqrt{\frac\pi2}(s_i^\top y)\operatorname{sign}(s_i^\top u).
\end{equation}
Then
\begin{equation}
    \operatorname{Var}(Z_i)
    \le \mathbb E Z_i^2
    =\frac\pi2\mathbb E(s_i^\top y)^2
    =\frac\pi2\|y\|_2^2.
\end{equation}
The average of $d$ independent terms has variance at most $\pi\|y\|_2^2/(2d)$.
\hfill$\square$

\paragraph{Non-unit residuals.}
If $r\ne 0$ has norm $\gamma=\|r\|_2$, write $u=r/\gamma$.
Then
\begin{equation}
    \operatorname{sign}(Sr)=\operatorname{sign}(Su),
\end{equation}
and the scaled QJL reconstruction
\begin{equation}
    \widetilde r_{\mathrm{qjl}}
    =\gamma\frac{\sqrt{\pi/2}}{d}S^\top\operatorname{sign}(Sr)
\end{equation}
satisfies
\begin{equation}
    \mathbb E\langle y,\widetilde r_{\mathrm{qjl}}\rangle
    =\langle y,r\rangle,
\end{equation}
with variance at most
\begin{equation}
    \frac{\pi}{2d}\gamma^2\|y\|_2^2.
\end{equation}

\subsubsection*{Inner-product TURBOQUANT}

The inner-product version combines two ideas:
\begin{enumerate}[leftmargin=2em]
    \item use MSE-optimized TURBOQUANT with $b-1$ bits per coordinate to get a good coarse approximation;
    \item use one additional bit per coordinate to apply QJL to the residual.
\end{enumerate}
The MSE stage makes the residual small.
The QJL stage makes the inner-product estimate unbiased.

\paragraph{Algorithm: \textsc{TURBOQUANT}$_{\mathrm{prod}}$.}
For a unit vector $x\in S^{d-1}$ and target bit-width $b\ge 1$:
\begin{enumerate}[leftmargin=2em]
    \item Run \textsc{TURBOQUANT}$_{\mathrm{mse}}$ at bit-width $b-1$:
    \[
        \mathrm{idx}=Q_{\mathrm{mse}}(x),
        \qquad
        \widetilde x_{\mathrm{mse}}=Q_{\mathrm{mse}}^{-1}(\mathrm{idx}).
    \]
    \item Form the residual
    \[
        r=x-\widetilde x_{\mathrm{mse}},
        \qquad
        \gamma=\|r\|_2.
    \]
    \item Store one QJL bit per coordinate:
    \[
        q=\operatorname{sign}(Sr)\in\{-1,+1\}^{d}.
    \]
    \item The compressed representation is
    \[
        (\mathrm{idx},q,\gamma).
    \]
    The index vector uses $(b-1)d$ bits, the QJL sign vector uses $d$ bits, so the main bit budget is $bd$ bits.
    \item Dequantize by
    \[
        \widetilde x_{\mathrm{qjl}}
        =\gamma\frac{\sqrt{\pi/2}}{d}S^\top q,
    \]
    and output
    \[
        \widetilde x=\widetilde x_{\mathrm{mse}}+\widetilde x_{\mathrm{qjl}}.
    \]
\end{enumerate}

\begin{theorem}[Inner-product guarantee for \textsc{TURBOQUANT}$_{\mathrm{prod}}$]
\label{thm:tq-prod}
For every unit vector $x\in S^{d-1}$ and every query vector $y\in\mathbb R^d$,
\begin{equation}
    \mathbb E\langle y,\widetilde x\rangle
    =\langle y,x\rangle.
\end{equation}
Moreover,
\begin{equation}
    D_{\mathrm{prod}}(Q_{\mathrm{prod}};x,y)
    \le
    \frac{\sqrt 3\,\pi^2\|y\|_2^2}{d}\,4^{-b}.
\end{equation}
For small bit-widths, the sharper values are approximately
\begin{equation}
    D_{\mathrm{prod}}
    \approx
    \frac{1.57}{d},
    \quad
    \frac{0.56}{d},
    \quad
    \frac{0.18}{d},
    \quad
    \frac{0.047}{d}
\end{equation}
for $b=1,2,3,4$, respectively, when $\|y\|_2=1$.
\end{theorem}

\paragraph{Proof.}
Condition on the MSE reconstruction $\widetilde x_{\mathrm{mse}}$.
Then the residual $r=x-\widetilde x_{\mathrm{mse}}$ is fixed for the QJL step.
By QJL unbiasedness,
\begin{align}
    \mathbb E[\langle y,\widetilde x\rangle\mid \widetilde x_{\mathrm{mse}}]
    &=\mathbb E[\langle y,\widetilde x_{\mathrm{mse}}+\widetilde x_{\mathrm{qjl}}\rangle\mid \widetilde x_{\mathrm{mse}}]\\
    &=\langle y,\widetilde x_{\mathrm{mse}}\rangle+\langle y,r\rangle\\
    &=\langle y,x\rangle.
\end{align}
Taking expectation over $\widetilde x_{\mathrm{mse}}$ proves unbiasedness.

For the distortion, again condition on $\widetilde x_{\mathrm{mse}}$:
\begin{align}
    \mathbb E\left[|\langle y,x\rangle-\langle y,\widetilde x\rangle|^2
    \mid \widetilde x_{\mathrm{mse}}\right]
    &=\mathbb E\left[|\langle y,r\rangle-\langle y,\widetilde x_{\mathrm{qjl}}\rangle|^2
    \mid \widetilde x_{\mathrm{mse}}\right]\\
    &=\operatorname{Var}(\langle y,\widetilde x_{\mathrm{qjl}}\rangle\mid \widetilde x_{\mathrm{mse}})\\
    &\le \frac{\pi}{2d}\|r\|_2^2\|y\|_2^2.
\end{align}
Now average over the MSE stage:
\begin{align}
    D_{\mathrm{prod}}
    &\le \frac{\pi}{2d}\|y\|_2^2\mathbb E\|x-\widetilde x_{\mathrm{mse}}\|_2^2.
\end{align}
The MSE stage uses $b-1$ bits per coordinate, so by Theorem~\ref{thm:tq-mse},
\begin{equation}
    \mathbb E\|x-\widetilde x_{\mathrm{mse}}\|_2^2
    \le
    \frac{\sqrt3\,\pi}{2}\,4^{-(b-1)}.
\end{equation}
Therefore
\begin{align}
    D_{\mathrm{prod}}
    &\le
    \frac{\pi}{2d}\|y\|_2^2
    \cdot
    \frac{\sqrt3\,\pi}{2}4^{-(b-1)}\\
    &=
    \frac{\sqrt3\,\pi^2\|y\|_2^2}{d}4^{-b}.
\end{align}
\hfill$\square$

\subsubsection*{Inner-product lower bound}

The MSE lower bound also implies an inner-product lower bound.
Let the reconstruction error be
\begin{equation}
    e=x-Q^{-1}(Q(x)).
\end{equation}
Then
\begin{equation}
    \|e\|_2^2
    =\sum_{j=1}^{d}\langle e_j,e\rangle^2,
\end{equation}
where $e_j$ is the $j$-th standard basis vector.
If every coordinate error had expected square less than $4^{-b}/d$, then the total MSE would be less than $4^{-b}$, contradicting the MSE lower bound.
Thus there exists some coordinate direction $e_j$ such that
\begin{equation}
    \mathbb E\left[
    |\langle e_j,x\rangle-\langle e_j,Q^{-1}(Q(x))\rangle|^2
    \right]
    \ge \frac1d4^{-b}.
\end{equation}
Equivalently, there exists a query direction $y$ with $\|y\|_2=1$ for which every quantizer suffers inner-product distortion at least
\begin{equation}
    D_{\mathrm{prod}}\ge \frac1d4^{-b}.
\end{equation}
For general $y$, scaling gives
\begin{equation}
    D_{\mathrm{prod}}\ge \frac{\|y\|_2^2}{d}4^{-b}
\end{equation}
for some hard direction with that norm.
Thus the inner-product version of TURBOQUANT also has the optimal bit-width exponent.

\subsubsection*{Why the two-stage design is natural}

The inner-product algorithm may look slightly unusual at first because it spends $b-1$ bits on MSE and one bit on a residual sketch.
The design is natural from the following decomposition:
\begin{equation}
    x=\widetilde x_{\mathrm{mse}}+r.
\end{equation}
Then
\begin{equation}
    \langle y,x\rangle
    =\langle y,\widetilde x_{\mathrm{mse}}\rangle+\langle y,r\rangle.
\end{equation}
The first term is computed from the MSE quantized vector.
The second term is estimated by QJL.
If the MSE stage is good, $\|r\|_2$ is small.
QJL has variance proportional to $\|r\|_2^2\|y\|_2^2/d$, so the residual sketch becomes accurate precisely because the first stage minimized residual norm.

\paragraph{Bias--variance tradeoff.}
The MSE stage alone has low variance but can be biased for inner products.
QJL alone is unbiased but has variance about $\pi\|y\|_2^2/(2d)$ for a unit vector.
The two-stage method keeps the unbiasedness of QJL while reducing its variance by applying it only to the residual.

\begin{center}
\begin{tabular}{c|c|c}
Method & Inner-product bias & Inner-product variance \\
\hline
MSE scalar quantizer after rotation & can be biased & small when $b$ is large \\
QJL only & zero & about $\pi\|y\|_2^2/(2d)$ \\
MSE $+$ QJL residual & zero & proportional to MSE residual
\end{tabular}
\end{center}

\subsubsection*{Online and data-oblivious nature}

A key practical feature of TURBOQUANT is that it is \emph{data-oblivious}.
The quantizer does not need to run $k$-means on a dataset of vectors.
It only needs:
\begin{enumerate}[leftmargin=2em]
    \item a public rotation $\Pi$;
    \item a precomputed scalar codebook for the sphere-coordinate density;
    \item for the inner-product version, a public Gaussian projection matrix $S$.
\end{enumerate}
Thus each vector can be quantized immediately as it arrives.
This is important for streaming settings such as KV cache quantization, where keys and values are generated online during inference.

\paragraph{Comparison with product quantization.}
Product quantization usually learns codebooks from a representative dataset.
This can work well for static vector databases, but it introduces preprocessing, indexing latency, and dependence on the training distribution.
TURBOQUANT instead uses a universal random rotation and scalar codebook.
Its strength is not that it learns the best possible empirical codebook for one dataset, but that it gives a fast, vectorizable, online quantizer with near-optimal worst-case distortion scaling.

\paragraph{Rotation cost.}
The mathematical description uses a Haar random orthogonal matrix, which can be generated by QR decomposition of a Gaussian matrix.
A dense rotation costs $O(d^2)$ operations per vector.
In practice, one can use structured rotations, block rotations, or hardware-friendly transforms to reduce cost.
The lecture-level theory should be read as explaining the geometry: the rotation spreads energy and makes coordinates look like a known scalar source.
Implementation choices decide how to approximate this geometry efficiently.

\subsubsection*{Application: KV cache quantization}

In a decoder-only transformer, the model stores key and value vectors for previous tokens.
For long contexts, this KV cache can dominate memory usage and memory bandwidth.
The attention score for a new query $q$ involves inner products
\begin{equation}
    \langle q,k_t\rangle
\end{equation}
with past keys $k_t$.
If the stored keys are quantized, the attention logits are perturbed.
Thus inner-product preservation is directly relevant.

TURBOQUANT is well matched to this setting for three reasons:
\begin{enumerate}[leftmargin=2em]
    \item \textbf{Online compression:} each new key/value vector can be compressed as soon as it is produced.
    \item \textbf{No calibration:} the codebook does not need to be trained on a calibration set.
    \item \textbf{Inner-product control:} the product version gives an unbiased inner-product estimator with variance scaling like $1/d$ times the residual energy.
\end{enumerate}

Empirically, the paper reports strong long-context performance at low bit-widths, including near quality-neutral results around $3.5$ bits per channel and modest degradation around $2.5$ bits per channel on long-context benchmarks.
The theoretical results do not by themselves prove downstream language-model quality, but they explain why geometric quantities used by attention can survive aggressive quantization.

\subsubsection*{Application: nearest-neighbor search}

In nearest-neighbor search, one stores a database of vectors $x_1,\dots,x_N$ and receives a query $y$.
The goal is to find vectors with large $\langle y,x_i\rangle$ or small Euclidean distance.
A compressed database is useful only if it preserves these scores well.

For unit-norm database vectors, the inner-product guarantee says
\begin{equation}
    \mathbb E\left[
        \left|\langle y,x_i\rangle-
        \langle y,\widetilde x_i\rangle\right|^2
    \right]
    \lesssim
    \frac{\|y\|_2^2}{d}4^{-b}.
\end{equation}
The $1/d$ factor is important.
In high dimension, random-projection inner-product errors average out across many independent rows of $S$.
Thus the same geometry that makes random rotations useful for MSE also helps inner-product estimation.

\subsubsection*{A compact summary of the proof architecture}

The complete TURBOQUANT analysis can be summarized as follows.

\begin{enumerate}[leftmargin=2em]
    \item \textbf{Random rotation creates a known source.}
    For any fixed $x\in S^{d-1}$, $\Pi x$ is uniform on the sphere, and each coordinate has the Beta-type density $f_d$.

    \item \textbf{High dimension makes the source nearly scalar Gaussian.}
    Each coordinate is approximately $\mathcal N(0,1/d)$, and small collections of coordinates are nearly independent.

    \item \textbf{Optimal scalar quantization gives vector MSE.}
    Quantizing each rotated coordinate with a Lloyd--Max quantizer gives total distortion $dC(f_d,b)$.

    \item \textbf{High-resolution scalar quantization gives the $4^{-b}$ exponent.}
    The Panter--Dite formula yields
    \[
        dC(f_d,b)\le \frac{\sqrt 3\pi}{2}4^{-b}.
    \]

    \item \textbf{Shannon's lower bound proves the exponent is unavoidable.}
    Any $bd$-bit vector quantizer has a hard input with MSE at least $4^{-b}$.

    \item \textbf{MSE alone is not enough for unbiased inner products.}
    At one bit, MSE quantization shrinks inner products by about $2/\pi$.

    \item \textbf{QJL corrects the bias on the residual.}
    MSE quantization first makes the residual small; QJL then estimates residual inner products without bias.
\end{enumerate}

\subsubsection*{Important caveats}

\paragraph{Expected distortion.}
The theorems are stated in expectation over the quantizer randomness.
In practice, a public random seed fixes the rotation and projection matrices.
For high dimension and many vectors, concentration usually makes a single draw behave close to the expectation, but the cleanest statement is an expected-distortion statement.

\paragraph{Side information.}
The algorithms store a vector norm, and the inner-product version stores a residual norm.
These are scalar side-information values.
The theoretical bit-width $b$ counts the dominant coordinatewise bits; the scalar overhead is negligible when $d$ is large but should be accounted for in an implementation.

\paragraph{Worst-case versus data-dependent optimality.}
TURBOQUANT is designed for worst-case vectors with a randomized, data-oblivious quantizer.
A data-dependent offline method may do better on a particular fixed dataset.
The point of the theorem is different: without making assumptions on the input distribution, TURBOQUANT matches the information-theoretic bit-width exponent up to constants.

\paragraph{From geometry to task quality.}
Preserving MSE and inner products is a strong geometric guarantee.
Downstream performance in language models or retrieval systems also depends on nonlinear model behavior, distribution shift, and implementation details.
The geometric theory explains why the method should be competitive, but empirical validation is still necessary for each application.

\subsubsection*{Takeaway}

TURBOQUANT is a clean example of how classical rate--distortion ideas become an algorithmic design principle.
The Gaussian rate--distortion function predicts the $4^{-b}$ scaling.
A random rotation turns a worst-case vector into a spherical random vector whose coordinates look Gaussian.
Scalar Lloyd--Max quantization after rotation achieves the right MSE scaling.
Finally, because MSE-optimized quantization is biased for inner products, a one-bit unbiased QJL sketch of the residual gives an inner-product-preserving version.

In one sentence:
\begin{center}
\emph{Rotate to Gaussianize, scalar-quantize to minimize MSE, and sketch the residual to preserve inner products.}
\end{center}
